\documentclass[a4paper]{article}

\usepackage{INTERSPEECH2021}
\usepackage[T1]{fontenc}
\usepackage{url}

\title{Joint Intent Classification and Slot Filling on ATIS:\\from Recurrent Networks to BERT}
\name{Andrea Giampietro}
\address{University of Trento}
\email{}

\begin{document}

\maketitle

\begin{abstract}
We train a single model to recognise the intent of a flight-domain request and to label its words with semantic
slots. On ATIS, a bidirectional LSTM with dropout raises the slot F1 of an LSTM baseline from 92.24 to 94.59 and
its intent accuracy from 93.44\% to 94.94\%. Fine-tuning BERT, with the slot labels aligned to the first WordPiece
of every word, reaches 95.49 slot F1 and 97.54\% intent accuracy (means over 5 runs).
\end{abstract}

\section{Introduction}

Spoken language understanding maps a request to a semantic frame: the intent says what the user wants and the slots
carry its arguments. In \emph{list flights from la guardia to burbank} the intent is \texttt{flight}, \emph{la
guardia} fills \texttt{fromloc.airport\_name} and \emph{burbank} fills \texttt{toloc.city\_name}. The two tasks are
related and are commonly learnt by one model with two outputs \cite{liu2016joint}. We study them on ATIS
\cite{hemphill1990atis}: part~A extends an LSTM baseline with a bidirectional encoder \cite{schuster1997birnn} and dropout \cite{srivastava2014dropout}, and part~B fine-tunes BERT
\cite{devlin2019bert} on both tasks at once \cite{chen2019jointbert}. The code is available at
\url{https://github.com/giampyhatesyou/nlu-ablations}.

\section{Implementation details}

\textbf{Data and evaluation.} We use a preprocessed version of ATIS with 4,978 training and 893 test utterances,
26 intents and 129 slot labels in IOB format \cite{ramshaw1995iob}. 10\% of the training set, stratified by intent,
is held out as dev set. Intents are evaluated by accuracy and slots by chunk-level F1 with the conlleval script
\cite{tjongkimsang2000conll}.

\textbf{Part A.} Words are embedded in 300 dimensions and encoded by an LSTM with 200 units per direction, which
skips the padding through packed sequences. A linear slot layer reads every encoder state and a linear intent layer
the final state; in the bidirectional model the intent layer receives the last forward state concatenated with the
backward state at the first word, which together summarise the whole utterance. Dropout ($p=0.3$) is applied to the
embeddings and before both output layers. The loss is the sum of the two cross-entropies, minimised with Adam
\cite{kingma2015adam} (lr $10^{-3}$, batch size 128, gradients clipped to norm 5).

\textbf{Part B.} \texttt{bert-base-uncased} feeds a linear intent layer with the pooled \texttt{[CLS]}
representation and a linear slot layer with the last-layer representation of every sub-token (dropout 0.1).
WordPiece \cite{wu2016wordpiece} splits 2.9\% of the test words, but 6.1\% of the words inside a slot (city and
airport names, aircraft and fare codes), and one test utterance in four contains a split word. Since ATIS labels
whole words, the label of a word is assigned to its first sub-token, while the other sub-tokens, \texttt{[CLS]},
\texttt{[SEP]} and the padding get the ignore index of the loss (Table~\ref{tab:align}). At inference the
prediction of a word is read from its first sub-token, so the slot F1 is computed on the original words with the
same script as part~A. The pre-trained encoder is loaded once per run and only the two output layers start from
scratch. We fine-tune with AdamW \cite{loshchilov2019adamw} (lr $5\cdot10^{-5}$, weight decay 0.01 except on biases
and LayerNorm), a linear warm-up over the first 10\% of the steps followed by a linear decay, batch size 32 and at
most 10 epochs.

\textbf{Model selection.} Each configuration is trained 5 times with different seeds. In every run the dev slot F1
is computed after each epoch, training stops after 10 (part~A) or 3 (part~B) epochs without improvement, and the
best weights are restored before testing. We report mean and sample standard deviation over the runs and compare
consecutive models with Welch's t-test ($\alpha = 0.05$).

\section{Results}

Table~\ref{tab:results} reports mean and sample standard deviation of the test scores over the five runs. The
bidirectional encoder gives the largest gain on slot filling (+1.89 F1, $p<0.001$), since the backward direction
lets the label of a word depend on the words that follow it. Dropout adds +0.46 F1 and +0.40 points of accuracy,
but with five runs neither gain is significant ($p=0.10$ and $p=0.06$): they are within the variation between
seeds. BERT is the best model on both tasks ($p \le 0.001$): the margin is large on intents (+2.60 points), where a
pre-trained representation of the whole sentence helps most, and smaller on slots (+0.90 F1), which in ATIS
largely depend on local lexical cues already captured by the BiLSTM; it also has the smallest spread on slots
($\pm$0.11 F1). Averaging over runs matters: the BERT run with the best dev F1, the one kept as
checkpoint, has the lowest test F1 of the five (95.33). Thanks to the first-sub-token alignment, the WordPiece
segmentation does not penalise BERT: all the models are scored on exactly the same words.

\begin{table}[t]
  \caption{Test results on ATIS: mean $\pm$ sample standard deviation over 5 runs, each evaluated at its best dev
  epoch. $^{*}$: significantly better than the previous row (Welch's t-test, $p<0.05$).}
  \label{tab:results}
  \centering
  % Generated by reports/analysis.py from results/: do not edit by hand
\begin{tabular}{l c c}
  \toprule
  Model & Slot F1 & Intent acc. (\%) \\
  \midrule
  LSTM & 92.24 $\pm$ 0.54 & 93.44 $\pm$ 0.47 \\
  BiLSTM & 94.13 $\pm$ 0.45$^{*}$ & 94.54 $\pm$ 0.15$^{*}$ \\
  BiLSTM + dropout & 94.59 $\pm$ 0.30 & 94.94 $\pm$ 0.35 \\
  BERT & \textbf{95.49} $\pm$ 0.11$^{*}$ & \textbf{97.54} $\pm$ 0.24$^{*}$ \\
  \bottomrule
\end{tabular}

\end{table}

\begin{table}[t]
  \caption{Training targets of BERT for the test utterance \emph{list flights from la guardia to burbank}: only the
  first sub-token of each word has a label.}
  \label{tab:align}
  \centering
  \small
  \begin{tabular}{l l l}
    \toprule
    Sub-token & Word & Target \\
    \midrule
    \texttt{[CLS]}    & --      & ignored \\
    \texttt{list}     & list    & \texttt{O} \\
    \texttt{flights}  & flights & \texttt{O} \\
    \texttt{from}     & from    & \texttt{O} \\
    \texttt{la}       & la      & \texttt{B-fromloc.airport\_name} \\
    \texttt{guard}    & guardia & \texttt{I-fromloc.airport\_name} \\
    \texttt{\#\#ia}   & guardia & ignored \\
    \texttt{to}       & to      & \texttt{O} \\
    \texttt{bu}       & burbank & \texttt{B-toloc.city\_name} \\
    \texttt{\#\#rba}  & burbank & ignored \\
    \texttt{\#\#nk}   & burbank & ignored \\
    \texttt{[SEP]}    & --      & ignored \\
    \bottomrule
  \end{tabular}
\end{table}

\bibliographystyle{IEEEtran}
\bibliography{references}

\end{document}
